\section{Proposed method}\label{sec:method}
AnytimeSwarm-CC takes three decisions in every replanning cycle (Fig.~\ref{fig:overview}, Algorithm~\ref{alg:cycle}).
The chance-constrained planning model decides which plans are acceptable under the calibrated execution error, the
stopping rule decides when to stop computing and publish, and the trigger decides when to replan. A feasibility-first
repair and an anytime ALNS produce the plan. Two earlier variants serve as baselines. AnytimeSwarm uses the same
stopping rule with a deterministic delay cost on the nominal model, and AnytimeSwarm-R uses an execution-aware delay
cost on the nominal model.

\begin{figure*}[t]
\centering
\fig{fig_overview}{\linewidth}
\caption{Overview of AnytimeSwarm-CC. Top: the replanning loop between the mission, the planning snapshot, the
replanner and the execution tier. Bottom: one replanning call. The calibrated error model gives a margin that grows
with the lead time, the chance constraint turns it into an effective deadline per task, the nominal fallback keeps the
tasks that no vehicle could otherwise reach, and the anytime search stops when the predicted gain no longer pays for the
delay. The plan is released after the emulated planning delay $\kappa\tau$ and the latency $L$.}\label{fig:overview}
\end{figure*}

\begin{algorithm}[t]
\caption{One replanning call of AnytimeSwarm-CC}\label{alg:cycle}
\small
\begin{algorithmic}[1]
\Require snapshot at time $t_0$: committed tasks, open tasks with $(r_j,d_j,w_j)$, UAV states; previous plan $\pi^-$;
error model $(b_0,b_1,s_0,s_1)$, safety factor $z$, expected latency $\bar L$, latency factor $\kappa$
\State $\tilde d_j \gets \min\{d_j,\; t_0+(d_j-t_0-b_0-z s_0)/(1+b_1+z s_1)\}$ for every open task
\Comment{Eq.~\eqref{eq:cc}}
\State $d'_j \gets d_j$ if no UAV can finish $j$ by $\tilde d_j$ but one can by $d_j$, else $d'_j\gets\tilde d_j$
\Comment{Eq.~\eqref{eq:fallback}}
\State $\pi \gets$ \Call{FeasibilityFirstRepair}{$\pi^-$, $d'$};\; $\pi^\ast \gets \pi$;\; $\tau \gets 0$
\Repeat
  \State destroy and repair $\pi$ with adaptively chosen operators; accept by simulated annealing
  \State \textbf{if} $J(\pi) > J(\pi^\ast)$ \textbf{then} $\pi^\ast \gets \pi$
  \State every 5\,ms: $\widehat{\Delta J} \gets$ learned prediction of the gain in the next $H$\,ms
\Until{$\widehat{\Delta J}\le C_{\mathrm{R}}(\tau)+\mu W H$ \textbf{or} $\tau\ge B$} \Comment{Eq.~\eqref{eq:stop}}
\State publish $\pi^\ast$; it reaches the UAVs at $t_0+\kappa\tau+L$
\State next call at the next event or $T_{\max}$ after $t_0$, whichever comes first
\end{algorithmic}
\end{algorithm}

\subsection{Repair and anytime search}
The previous plan is restricted to the open tasks and walked once, keeping a task only if the route up to it is still
feasible, and tasks of failed UAVs and new tasks are inserted by parallel regret-2
insertion~\cite{potvin1993parallel,ropke2006adaptive}. This feasibility-first repair yields a publishable plan within
milliseconds. ALNS~\cite{ropke2006adaptive} then improves it with four destroy operators (random, worst,
related~\cite{shaw1998cp} and route-segment removal), four greedy or regret-2 repair operators with and without noise,
adaptive operator selection and simulated-annealing acceptance~\cite{kirkpatrick1983annealing}. All intermediate
solutions are feasible, so the best plan found so far can be published at any time. The temperature decreases per
iteration and all variants share the random seed, so every stopping rule follows the same search trajectory and differs
only in the moment at which it stops.

\subsection{Calibrated finish-time error}
The error of task $j$ is modelled as $e_j\sim\mathcal{N}(m(h_j),s(h_j)^2)$, where $h_j=f_j-t_0$ is its lead time, the
nominal finish time minus the time $t_0$ of the planning snapshot, and
\begin{equation}\label{eq:errmodel}
m(h)=b_0+b_1h,\qquad s(h)=s_0+s_1h ,
\end{equation}
because errors accumulate along a route. The parameters are fitted per disturbance level on development missions, the
mean by least squares and the spread by least squares on the absolute residuals scaled by $\sqrt{\pi/2}$, which is less
sensitive to the heavy tail of the errors. Under the mid and high disturbances of the main studies, both the mean and
the spread grow with the lead time (Supplementary Table~S2). The nominal model is therefore optimistic, and
increasingly so along a route.

\subsection{Chance-constrained planning model}\label{sec:cc}
The proposed model plans for the calibrated execution. A planned finish time $f_j$ is accepted only if
\begin{equation}\label{eq:ccfinish}
f_j+m(f_j-t_0)+z\,s(f_j-t_0)\;\le\;d_j ,
\end{equation}
where $z\ge0$ is a safety factor. The margin $m+zs$ depends on the finish time itself, so it is small for tasks served
soon and large for tasks served late in a route. Because $m$ and $s$ are linear, the constraint reduces to an effective
deadline,
\begin{equation}\label{eq:cc}
f_j\le\tilde d_j,\qquad \tilde d_j=t_0+\frac{d_j-t_0-b_0-z s_0}{1+b_1+z s_1}.
\end{equation}
The repair and the search use $\tilde d_j$ instead of $d_j$, while the executor and the reward still use the true
deadline.

\begin{prop}\label{prop:cc}
Let $1+b_1+z s_1>0$. (i) A finish time satisfies Eq.~\eqref{eq:ccfinish} if and only if $f_j\le\tilde d_j$. (ii) Under
the error model of Eq.~\eqref{eq:errmodel}, a task planned to finish at $f_j\le\tilde d_j$ meets its deadline with
probability at least $\Phi(z)$ if its route starts as planned. (iii) With $b_0=b_1=s_0=s_1=0$, $\tilde d_j=d_j$ and the
planning model is the nominal one.
\end{prop}
\begin{proof}
Substituting Eq.~\eqref{eq:errmodel} into Eq.~\eqref{eq:ccfinish} gives
$f_j(1+b_1+zs_1)\le d_j-b_0-zs_0+(b_1+zs_1)t_0$, and dividing by the positive factor gives (i). For (ii), the actual
finish is $F_j=f_j+e_j$ with $e_j\sim\mathcal{N}(m(h_j),s(h_j)^2)$ and $h_j=f_j-t_0$. By (i),
$d_j-f_j-m(h_j)\ge z\,s(h_j)$, hence $\Pr[F_j\le d_j]=\Phi\big((d_j-f_j-m(h_j))/s(h_j)\big)\ge\Phi(z)$. Part (iii) is
immediate.
\end{proof}

The effective deadline turns a chance constraint on every task into a deterministic constraint of the same form as the
original one. The insertion heuristics and the feasibility checks of the search are therefore reused unchanged, and
computing $\tilde d_j$ costs one division per task. With $z=0$, the model corrects the mean error only and plans for the
median finish time. Larger values of $z$ protect more tasks against lateness but leave fewer tasks feasible, so $z$
trades the risk of a late service against the reward of attempting a task. It is chosen per disturbance level on
development missions.

\subsection{Nominal fallback for given-up tasks}\label{sec:fallback}
A hard chance constraint has a side effect when the time windows are short compared with the spread of the error. If no
vehicle can finish task $j$ before $\tilde d_j$, even on a direct flight, the task is infeasible for the planner and is
never attempted. A vehicle may still serve it in time with a probability below $\Phi(z)$ but well above zero, and it may
have nothing better to do. The fallback keeps such tasks. Let $\tilde H$ be the set of tasks that no vehicle can finish
by $\tilde d_j$ on a direct flight, and $H$ the set of tasks that no vehicle can finish by $d_j$. The planning deadline is
\begin{equation}\label{eq:fallback}
d'_j=\begin{cases} d_j, & j\in\tilde H\setminus H,\\ \tilde d_j, & \text{otherwise.}\end{cases}
\end{equation}
The constraints of all other tasks are unchanged, so every other task that is served keeps the guarantee of
Proposition~\ref{prop:cc}, while a fallback task carries none. Finding $\tilde H$ and $H$ requires one direct-flight
check per task and vehicle, which is negligible compared with the search. The fallback was designed after the held-out
results of study~2 and is evaluated on new seeds in study~3.

\subsection{Stopping rule and trigger}\label{sec:stop}
Every 5\,ms, the stopping rule compares the predicted gain $\widehat{\Delta J}$ of searching $H=10$\,ms longer with the
cost of doing so, and stops at time $\tau$ if
\begin{equation}\label{eq:stop}
\widehat{\Delta J}\;\le\;C_{\mathrm{R}}(\tau)+\mu W H ,
\end{equation}
with a price of computation $\mu=10^{-5}$ per millisecond. The gain is predicted from eleven features of the search
history by a Poisson generalised linear model~\cite{mccullagh1989glm} fitted on development trajectories. The delay
cost $C_{\mathrm{R}}(\tau)=\Lambda_{\mathrm{R}}(\tau+H)-\Lambda_{\mathrm{R}}(\tau)$ is the increase of the expected value
of the tasks that a later start would lose,
\begin{equation}\label{eq:risk}
\Lambda_{\mathrm{R}}(\tau)=\sum_j v_j\,\Phi\!\left(\frac{\sigma_j(\tau)+\bar L+m(h_j)-\mathrm{tol}_j}{s(h_j)}\right),
\end{equation}
where $\sigma_j(\tau)$ is the start shift of the route of task $j$ caused by waiting for the plan, $\mathrm{tol}_j$ the
margin of the task to its deadline plus the waiting times before it on its route~\cite{savelsbergh1992vrptw}, $\bar L$
the expected latency and $v_j$ the value of the task in Eq.~\eqref{eq:obj}. On the effective deadlines, the mean error
is already part of the tolerance, so AnytimeSwarm-CC uses $m\equiv0$. $\Lambda_{\mathrm{R}}$ does not decrease with
$\tau$, and it reduces to the exact loss of a deterministic plan when the spread vanishes (Supplementary
Section~S1). A hard budget $B$ bounds every call. The replanner replans on every event and at least every
$T_{\max}=60$\,s of mission time, a rolling-horizon policy in the terminology of
rescheduling~\cite{vieira2003rescheduling,ouelhadj2009survey}.
